\documentclass[10pt,a4paper]{article}
\usepackage[T1]{fontenc}
\usepackage{lmodern}
\usepackage{amsmath,amssymb,amsthm,mathtools}
\usepackage[margin=24mm]{geometry}
\usepackage{microtype}
\usepackage{booktabs,array,longtable}
\usepackage{xcolor}
\usepackage{hyperref}
\usepackage{fancyhdr}
\hypersetup{colorlinks=true,linkcolor=black,citecolor=black,urlcolor=blue,pdfauthor={Matthew J. Colbrook}}
\pagestyle{fancy}
\fancyhf{}
\fancyhead[L]{\small Resolution / MI-26}
\fancyhead[R]{\small 11 September 2026}
\fancyfoot[C]{\thepage}
\setlength{\headheight}{14pt}
\setlength{\parskip}{4pt}
\setlength{\parindent}{0pt}
\newtheorem{theorem}{Theorem}[section]
\newtheorem{proposition}[theorem]{Proposition}
\newtheorem{lemma}[theorem]{Lemma}
\newtheorem{corollary}[theorem]{Corollary}
\theoremstyle{remark}
\newtheorem{remark}[theorem]{Remark}
\newcommand{\M}{\mathbb M}
\newcommand{\C}{\mathbb C}
\newcommand{\R}{\mathbb R}
\newcommand{\tr}{\operatorname{tr}}
\newcommand{\diag}{\operatorname{diag}}
\newcommand{\spec}{\operatorname{spec}}
\newcommand{\rank}{\operatorname{rank}}
\newcommand{\abs}[1]{\lvert #1\rvert}
\newcommand{\norm}[1]{\lVert #1\rVert}
\newcommand{\opnorm}[1]{\lVert #1\rVert_{\infty}}
\newcommand{\schatten}[2]{\lVert #1\rVert_{#2}}
\newcommand{\symmod}{\mathcal S}
\newcommand{\maxmod}{\mathcal M}
\newcommand{\draftnotice}{\begin{center}\small\textit{Proposed mathematical resolution. Complete proof supplied; not submitted or peer reviewed.}\end{center}}

\usepackage{xurl}
\setlength{\emergencystretch}{3em}
\allowdisplaybreaks[1]

\title{MI-26: Resolution}
\author{Matthew J. Colbrook\\
\small Department of Applied Mathematics and Theoretical Physics\\
\small University of Cambridge, Cambridge, United Kingdom\\
\small\href{mailto:m.colbrook@damtp.cam.ac.uk}{m.colbrook@damtp.cam.ac.uk}}
\date{11 September 2026}
\begin{document}
\maketitle
\textbf{Repository status: Solved.}
The complete argument received an independent Codex-agent PASS review
against the exact catalog target.
The original draft was AI-assisted; the reviewed proof below is unchanged.
This is independent agent verification, not external human peer review or
formal proof certification. \href{https://github.com/ajt60gaibb/OpenProblemsInNLA/blob/main/references/colbrook-matrix-2026-09-11/verification/reviews/MI-26-review.md}{Detailed review and scope}.
\par\medskip
% BEGIN REVIEWED PROOF
\section{MI-26: concavity alone does not give two-unitary subadditivity}
\label{sec:mi26}

The target asks whether, for every real-valued concave function $f$ on $[0,\infty)$ with $f(0)\geq0$ and every positive semidefinite $A,B$, there are unitaries $U,V$ such that
\begin{equation}\label{eq:mi26-target}
 f(A+B)\leq Uf(A)U^*+Vf(B)V^*.
\end{equation}
It is explicitly the removal of monotonicity that is at issue \cite[Problem 5]{AK2012}; see also \cite[Remark 3.13]{BL2011}.

\begin{theorem}\label{thm:mi26}
The assertion \eqref{eq:mi26-target} is false in dimension two. A counterexample is provided by the polynomial $f(t)=t-t^2$ and two real rank-one projections. It remains false when both inputs are required to be positive definite.
\end{theorem}

\begin{proof}
The function $f(t)=t-t^2$ is concave on $[0,\infty)$ and satisfies $f(0)=0$. Let
\[
 P=\begin{pmatrix}1&0\\0&0\end{pmatrix},\qquad
 Q=\frac1{25}\begin{pmatrix}9&12\\12&16\end{pmatrix}.
\]
Both matrices are orthogonal projections, so $f(P)=f(Q)=0$. However,
\[
 f(P+Q)=P+Q-(P+Q)^2=-(PQ+QP)
 =\frac1{25}\begin{pmatrix}-18&-12\\-12&0\end{pmatrix}.
\]
Its eigenvalues are $6/25$ and $-24/25$; in particular $(1,-2)^T$ is an eigenvector for $6/25$. Thus $f(P+Q)\not\leq0$, whereas the right-hand side of \eqref{eq:mi26-target} is zero for every $U,V$.

For positive definite inputs, take
\[
 A=P+\frac1{20}I_2,\qquad B=Q+\frac1{20}I_2.
\]
Each has eigenvalues $21/20$ and $1/20$, and consequently
\[
 \lambda_{\max}(f(A))=\lambda_{\max}(f(B))=\frac{19}{400}.
\]
The eigenvalues of $P+Q$ are $8/5$ and $2/5$. Thus $A+B$ has eigenvalues $17/10$ and $1/2$, giving
\[
 \lambda_{\max}(f(A+B))=\frac14.
\]
For every pair of unitaries,
\[
 Uf(A)U^*+Vf(B)V^*\leq\frac{19}{200}I_2.
\]
Since $1/4>19/200$, domination \eqref{eq:mi26-target} is again impossible.
\end{proof}

\begin{remark}[The function class matters]
This counterexample uses a concave function that is real-valued, not nonnegative on the whole half-line. It does not contradict the theorem for nonnegative concave functions. In fact, a concave function that is nonnegative everywhere on $[0,\infty)$ must be nondecreasing: a negative secant slope, by concavity, would force the function eventually to become negative. The repository target and the cited open question allow real-valued concave functions, so $f(t)=t-t^2$ belongs to the stated class.
\end{remark}

% END REVIEWED PROOF
\begin{thebibliography}{1}

\bibitem{AK2012}
Koenraad M.~R. Audenaert and Fuad Kittaneh.
\newblock Problems and conjectures in matrix and operator inequalities.
\newblock arXiv:1201.5232, 2012.
\newblock Version consulted: v3; Problem 5, printed page 3, and Problem 7,
  printed pages 12--13.
\newblock URL: \url{https://arxiv.org/abs/1201.5232}.

\bibitem{BL2011}
Jean-Christophe Bourin and Eun-Young Lee.
\newblock Unitary orbits of {Hermitian} operators with convex or concave
  functions.
\newblock arXiv:1109.2384, 2011.
\newblock Remark 3.13.
\newblock URL: \url{https://arxiv.org/abs/1109.2384}.

\end{thebibliography}

\end{document}
